\chapter{PNML and JSON Data I/O}
\label{chap:data-io}

This chapter covers the application's data persistence capabilities: how Petri nets are saved to files and loaded from files. The I/O services act as translators between external file formats and the internal data structures.

\section{Problem Statement}
\label{sec:io-problem}

After spending time building a Petri net, users need to:
\begin{itemize}
    \item \textbf{Save} their work for later continuation
    \item \textbf{Share} models with colleagues
    \item \textbf{Load} existing models from other tools
\end{itemize}

Without I/O capabilities, all work would be lost when closing the browser. The application supports two file formats to address different needs.

\section{Supported Formats}
\label{sec:io-formats}

\begin{table}[htb]
    \centering
    \begin{tabular}{lll}
        \toprule
        \textbf{Feature} & \textbf{PNML} & \textbf{Custom JSON} \\
        \midrule
        Format type & XML-based & JSON \\
        Standard & ISO/IEC 15909-2 & Application-specific \\
        Complexity & Verbose, detailed & Concise, simple \\
        Interoperability & WoPeD, CPN Tools, etc. & This application only \\
        Human-readable & Moderate & High \\
        \bottomrule
    \end{tabular}
    \caption{Comparison of supported file formats}
    \label{tab:format-comparison}
\end{table}

\section{Custom JSON Format}
\label{sec:json-format}

The custom JSON format is designed for simplicity and ease of parsing:

\begin{lstlisting}[language=TypeScript, caption={JSON format interface (src/app/classes/json-petri-net.ts)}]
/**
 * Structure of our custom JSON format for Petri nets.
 */
export interface JsonPetriNet {
    /** List of place IDs, e.g., ["p1", "p2", "p3"] */
    places: string[];
    
    /** List of transition IDs, e.g., ["t1", "t2"] */
    transitions: string[];
    
    /** Arc definitions as "source,target": weight pairs */
    arcs?: { [idPair: string]: number };  // e.g., "p1,t1": 1
    
    /** Unique action labels for transitions */
    actions?: string[];
    
    /** Transition labels as id: label pairs */
    labels?: { [transitionId: string]: string };
    
    /** Initial marking as id: tokenCount pairs */
    marking?: { [placeId: string]: number };
    
    /** Element positions as id: {x, y} pairs */
    layout?: { [id: string]: Coords | Coords[] };
}

export interface Coords {
    x: number;
    y: number;
}
\end{lstlisting}

\subsection{Example JSON File}

\begin{lstlisting}[language=json, caption={Example Petri net in JSON format}]
{
  "places": ["p1", "p2", "p3"],
  "transitions": ["t1", "t2"],
  "arcs": {
    "p1,t1": 1,
    "t1,p2": 1,
    "p2,t2": 1,
    "t2,p3": 1
  },
  "marking": {
    "p1": 1,
    "p2": 0,
    "p3": 0
  },
  "labels": {
    "t1": "Stamp Letter",
    "t2": "Mail Letter"
  },
  "layout": {
    "p1": {"x": 100, "y": 100},
    "t1": {"x": 200, "y": 100},
    "p2": {"x": 300, "y": 100},
    "t2": {"x": 400, "y": 100},
    "p3": {"x": 500, "y": 100}
  }
}
\end{lstlisting}

\section{Exporting to JSON}
\label{sec:json-export}

The \texttt{ExportJsonDataService} converts the internal data model to JSON format:

\begin{lstlisting}[language=TypeScript, caption={JSON export service (src/app/tr-services/export-json-data.service.ts)}]
import { Injectable } from '@angular/core';
import { DataService } from './data.service';
import { JsonPetriNet, Coords } from '../classes/json-petri-net';

@Injectable({ providedIn: 'root' })
export class ExportJsonDataService {
    constructor(private dataService: DataService) {}

    /**
     * Export the current Petri net as a JSON file download.
     */
    public exportAsJson(): void {
        // 1. Build the JsonPetriNet structure
        const jsonObj = this.buildJsonStructure();
        
        // 2. Serialize to formatted JSON string
        const jsonString = JSON.stringify(jsonObj, null, 2);
        
        // 3. Trigger browser download
        this.downloadFile(jsonString, 'petri-net.json', 'application/json');
    }

    private buildJsonStructure(): JsonPetriNet {
        const json: JsonPetriNet = {
            places: [],
            transitions: [],
            arcs: {},
            marking: {},
            labels: {},
            layout: {},
        };

        // Collect places
        for (const place of this.dataService.getPlaces()) {
            json.places.push(place.id);
            json.layout![place.id] = { x: place.position.x, y: place.position.y };
            
            if (place.token > 0) {
                json.marking![place.id] = place.token;
            }
        }

        // Collect transitions
        for (const transition of this.dataService.getTransitions()) {
            json.transitions.push(transition.id);
            json.layout![transition.id] = { 
                x: transition.position.x, 
                y: transition.position.y 
            };
            
            if (transition.label) {
                json.labels![transition.id] = transition.label;
            }
        }

        // Collect arcs
        for (const arc of this.dataService.getArcs()) {
            const key = `${arc.from.id},${arc.to.id}`;
            json.arcs![key] = arc.weight;
            
            // Store anchor points if present
            if (arc.anchors.length > 0) {
                json.layout![key] = arc.anchors.map(p => ({ x: p.x, y: p.y }));
            }
        }

        return json;
    }

    private downloadFile(content: string, filename: string, mimeType: string): void {
        const blob = new Blob([content], { type: mimeType });
        const url = URL.createObjectURL(blob);
        
        const link = document.createElement('a');
        link.href = url;
        link.download = filename;
        link.click();
        
        URL.revokeObjectURL(url);
    }
}
\end{lstlisting}

\section{Importing from JSON}
\label{sec:json-import}

The \texttt{ParserService} converts JSON files back to internal objects:

\begin{lstlisting}[language=TypeScript, caption={JSON parser service (src/app/tr-services/parser.service.ts)}]
import { Injectable } from '@angular/core';
import { JsonPetriNet, Coords } from '../classes/json-petri-net';
import { Place } from '../tr-classes/petri-net/place';
import { Transition } from '../tr-classes/petri-net/transition';
import { Arc } from '../tr-classes/petri-net/arc';
import { Point } from '../tr-classes/petri-net/point';

@Injectable({ providedIn: 'root' })
export class ParserService {
    /** Flag: true if some elements had no position data */
    incompleteLayoutData: boolean = false;

    /**
     * Parse a JSON string into Petri net elements.
     * Returns arrays of places, transitions, arcs, and actions.
     */
    parse(jsonText: string): [Place[], Transition[], Arc[], string[]] {
        this.incompleteLayoutData = false;
        
        // 1. Parse JSON string to object
        const data: JsonPetriNet = JSON.parse(jsonText);
        
        const places: Place[] = [];
        const transitions: Transition[] = [];
        const arcs: Arc[] = [];
        const actions: string[] = data.actions || [];

        // 2. Create Place objects
        for (const placeId of data.places) {
            const position = this.getPosition(data, placeId);
            const tokens = data.marking?.[placeId] || 0;
            const label = data.labels?.[placeId];
            
            places.push(new Place(tokens, position, placeId, label));
        }

        // 3. Create Transition objects
        for (const transitionId of data.transitions) {
            const position = this.getPosition(data, transitionId);
            const label = data.labels?.[transitionId];
            
            transitions.push(new Transition(position, transitionId, label));
        }

        // 4. Create Arc objects and link to nodes
        if (data.arcs) {
            for (const [idPair, weight] of Object.entries(data.arcs)) {
                const [fromId, toId] = idPair.split(',');
                
                const fromNode = this.findNode(fromId, places, transitions);
                const toNode = this.findNode(toId, places, transitions);
                
                if (fromNode && toNode) {
                    const anchors = this.getAnchors(data, idPair);
                    const arc = new Arc(fromNode, toNode, weight, anchors);
                    arcs.push(arc);
                    
                    // Register with transition
                    if (toNode instanceof Transition) {
                        toNode.appendPreArc(arc);
                    }
                    if (fromNode instanceof Transition) {
                        fromNode.appendPostArc(arc);
                    }
                }
            }
        }

        return [places, transitions, arcs, actions];
    }

    /**
     * Get position from layout data, or return default (0,0).
     */
    private getPosition(data: JsonPetriNet, id: string): Point {
        if (data.layout && data.layout[id]) {
            const coords = data.layout[id] as Coords;
            return new Point(coords.x, coords.y);
        }
        
        this.incompleteLayoutData = true;
        return new Point(0, 0);
    }

    /**
     * Get arc anchor points from layout data.
     */
    private getAnchors(data: JsonPetriNet, arcId: string): Point[] {
        if (data.layout && Array.isArray(data.layout[arcId])) {
            return (data.layout[arcId] as Coords[]).map(c => new Point(c.x, c.y));
        }
        return [];
    }

    /**
     * Find a node (place or transition) by ID.
     */
    private findNode(id: string, places: Place[], transitions: Transition[]): Place | Transition | undefined {
        return places.find(p => p.id === id) || transitions.find(t => t.id === id);
    }
}
\end{lstlisting}

\section{PNML Format}
\label{sec:pnml-format}

PNML (Petri Net Markup Language) is an XML-based international standard:

\begin{lstlisting}[language=XML, caption={Example PNML document structure}]
<?xml version="1.0" encoding="UTF-8"?>
<pnml>
  <net id="net1" type="http://www.pnml.org/version-2009/grammar/ptnet">
    <page id="page1">
      <place id="p1">
        <name><text>Unstamped Letter</text></name>
        <graphics>
          <position x="100" y="100"/>
        </graphics>
        <initialMarking>
          <text>1</text>
        </initialMarking>
      </place>
      
      <transition id="t1">
        <name><text>Stamp Letter</text></name>
        <graphics>
          <position x="200" y="100"/>
        </graphics>
      </transition>
      
      <arc id="a1" source="p1" target="t1">
        <inscription><text>1</text></inscription>
      </arc>
    </page>
  </net>
</pnml>
\end{lstlisting}

\section{PNML Import and Export}
\label{sec:pnml-service}

The \texttt{PnmlService} handles bidirectional PNML conversion:

\begin{lstlisting}[language=TypeScript, caption={PNML service (src/app/tr-services/pnml.service.ts) - Import}]
import { Injectable } from '@angular/core';
import { xml2js } from 'xml-js';
import { DataService } from './data.service';
import { Place } from '../tr-classes/petri-net/place';
import { Transition } from '../tr-classes/petri-net/transition';
import { Arc } from '../tr-classes/petri-net/arc';
import { Point } from '../tr-classes/petri-net/point';

@Injectable({ providedIn: 'root' })
export class PnmlService {
    incompleteLayoutData: boolean = false;

    constructor(
        private dataService: DataService,
        private layoutService: LayoutSugiyamaService
    ) {}

    /**
     * Parse PNML XML string and load into DataService.
     */
    parse(xmlString: string): [Place[], Transition[], Arc[], string[]] {
        this.incompleteLayoutData = false;
        
        // 1. Convert XML to JavaScript object
        const xmlObj = xml2js(xmlString, { compact: false });
        
        // 2. Find all place, transition, and arc elements
        const pnmlPlaces = this.findElements(xmlObj, 'place');
        const pnmlTransitions = this.findElements(xmlObj, 'transition');
        const pnmlArcs = this.findElements(xmlObj, 'arc');
        
        // 3. Convert to internal objects
        const places = this.parsePlaces(pnmlPlaces);
        const transitions = this.parseTransitions(pnmlTransitions);
        const arcs = this.parseArcs(pnmlArcs, places, transitions);
        const actions = this.extractActions(transitions);
        
        // 4. Update DataService
        this.dataService.places = places;
        this.dataService.transitions = transitions;
        this.dataService.arcs = arcs;
        this.dataService.actions = actions;
        
        // 5. Apply automatic layout if positions were missing
        if (this.incompleteLayoutData) {
            this.layoutService.applySugiyamaLayout();
        }
        
        this.dataService.triggerDataChanged(true);
        
        return [places, transitions, arcs, actions];
    }

    /**
     * Parse a PNML place element into a Place object.
     */
    private parsePlace(element: any): Place {
        const id = element.attributes.id;
        
        // Extract name from nested <name><text> elements
        const nameElement = this.findChild(element, 'name');
        const label = nameElement ? this.getTextContent(nameElement) : undefined;
        
        // Extract position from <graphics><position>
        const position = this.extractPosition(element);
        
        // Extract initial marking from <initialMarking><text>
        const markingElement = this.findChild(element, 'initialMarking');
        const tokens = markingElement ? parseInt(this.getTextContent(markingElement)) || 0 : 0;
        
        return new Place(tokens, position, id, label);
    }

    /**
     * Extract position from graphics element, or mark as incomplete.
     */
    private extractPosition(element: any): Point {
        const graphics = this.findChild(element, 'graphics');
        const position = graphics ? this.findChild(graphics, 'position') : null;
        
        if (position?.attributes?.x && position?.attributes?.y) {
            return new Point(
                parseFloat(position.attributes.x),
                parseFloat(position.attributes.y)
            );
        }
        
        this.incompleteLayoutData = true;
        return new Point(0, 0);
    }
    
    // ... additional helper methods ...
}
\end{lstlisting}

\subsection{PNML Export}

\begin{lstlisting}[language=TypeScript, caption={PNML export method}]
/**
 * Generate PNML XML string from current DataService content.
 */
public getPNML(): string {
    const places = this.dataService.getPlaces();
    const transitions = this.dataService.getTransitions();
    const arcs = this.dataService.getArcs();
    
    // Build XML structure
    let xml = `<?xml version="1.0" encoding="UTF-8"?>
<pnml>
  <net id="net1" type="http://www.pnml.org/version-2009/grammar/ptnet">
    <page id="page1">
`;
    
    // Add places
    for (const place of places) {
        xml += this.placeToXml(place);
    }
    
    // Add transitions
    for (const transition of transitions) {
        xml += this.transitionToXml(transition);
    }
    
    // Add arcs
    for (const arc of arcs) {
        xml += this.arcToXml(arc);
    }
    
    xml += `    </page>
  </net>
</pnml>`;
    
    return this.formatXml(xml);
}

private placeToXml(place: Place): string {
    return `      <place id="${place.id}">
        <name><text>${place.label || ''}</text></name>
        <graphics>
          <position x="${place.position.x}" y="${place.position.y}"/>
        </graphics>
        <initialMarking><text>${place.token}</text></initialMarking>
      </place>
`;
}

private transitionToXml(transition: Transition): string {
    return `      <transition id="${transition.id}">
        <name><text>${transition.label || ''}</text></name>
        <graphics>
          <position x="${transition.position.x}" y="${transition.position.y}"/>
        </graphics>
      </transition>
`;
}

private arcToXml(arc: Arc): string {
    return `      <arc id="arc-${arc.from.id}-${arc.to.id}" source="${arc.from.id}" target="${arc.to.id}">
        <inscription><text>${arc.weight}</text></inscription>
      </arc>
`;
}
\end{lstlisting}

\section{Import Flow Visualization}
\label{sec:io-flow}

\begin{figure}[htb]
    \centering
    \begin{tikzpicture}[
        node distance=1.5cm,
        box/.style={rectangle, draw, minimum width=2.5cm, minimum height=0.7cm, font=\small},
        decision/.style={diamond, draw, aspect=2, font=\small},
    ]
        \node[box] (file) {User selects file};
        \node[box, below=of file] (read) {Read file content};
        \node[decision, below=of read] (format) {Format?};
        \node[box, below left=of format] (json) {ParserService};
        \node[box, below right=of format] (pnml) {PnmlService};
        \node[box, below=2.5cm of format] (data) {Update DataService};
        \node[decision, below=of data] (layout) {Layout missing?};
        \node[box, below left=of layout] (sugiyama) {Apply Sugiyama};
        \node[box, below=2cm of layout] (notify) {triggerDataChanged()};
        \node[box, below=of notify] (redraw) {Canvas redraws};
        
        \draw[->] (file) -- (read);
        \draw[->] (read) -- (format);
        \draw[->] (format) -- node[above left] {.json} (json);
        \draw[->] (format) -- node[above right] {.pnml} (pnml);
        \draw[->] (json) |- (data);
        \draw[->] (pnml) |- (data);
        \draw[->] (data) -- (layout);
        \draw[->] (layout) -- node[above left] {yes} (sugiyama);
        \draw[->] (layout) -- node[right] {no} (notify);
        \draw[->] (sugiyama) |- (notify);
        \draw[->] (notify) -- (redraw);
    \end{tikzpicture}
    \caption{File import processing flow}
    \label{fig:import-flow}
\end{figure}

\section{Handling Missing Layout Data}
\label{sec:missing-layout}

When loading files without position information (common with PNML files from other tools), the services set \texttt{incompleteLayoutData = true}. The calling code then automatically applies the Sugiyama layout algorithm (covered in the next chapter) to generate readable positions.

\section{Summary}
\label{sec:io-summary}

The I/O services enable data persistence and interoperability:

\begin{itemize}
    \item \textbf{ExportJsonDataService}: Exports to custom JSON format
    \item \textbf{ParserService}: Imports from custom JSON format
    \item \textbf{PnmlService}: Bidirectional PNML (XML) conversion
    \item \textbf{Automatic layout}: Applied when position data is missing
    \item \textbf{Download trigger}: Browser's download API for file saving
\end{itemize}

The next chapter covers the \textbf{Layout Algorithms} that automatically arrange Petri net elements into readable configurations.
